\chapter{Introdução}
\label{cap:introducao}

\abbrev{HPC}{High-Performance Computing}
\abbrev{MPI}{Message Passing Interface}
\abbrev{E/S}{Entrada/Saída}
\abbrev{SDDP}{Stochastic Dual Dynamic Programming}

O aumento da complexidade e do volume de dados em aplicações científicas e
de engenharia tem tornado a computação paralela uma ferramenta indispensável
para atender às demandas de desempenho~\cite{Kang:20}. Entretanto,
operações de entrada e saída (E/S) continuam sendo um dos principais gargalos
em ambientes de alto desempenho (HPC), particularmente em aplicações que
processam dados massivos~\cite{Luu:15, Xie:20}. Nesse contexto, o uso de
bibliotecas como o MPI-IO, combinadas com sistemas de arquivos paralelos como
o Lustre, tornou-se uma abordagem fundamental para lidar com a escalabilidade
e a eficiência no acesso aos dados~\cite{mpi50Cite, LustreManual}.

Paralelamente, a crescente inserção de fontes renováveis na matriz
energética, como solar e eólica, tem introduzido novas dificuldades na
formulação e resolução de problemas de otimização associados à operação dos
sistemas elétricos. A variabilidade e incerteza dessas fontes requerem
modelos computacionais mais complexos e dependentes de grandes volumes de
dados, cuja execução eficiente depende tanto do paralelismo computacional
quanto de soluções eficazes de E/S~\cite{Ajc:10, Ssr:19}.

\textcolor{red}{\textbf{ToDp: precisamos expandir esse parágrafo, explicar o
que implica essa mudança de paradigma na simulação tanto para o uso de
recursos computacionais como para o uso dos resultados pelos usuários do
SDDP.}}

Assim como outras aplicações baseadas em tarefas independentes, o modelo
SDDP (\textit{Stochastic Dual Dynamic Programming}) é amplamente utilizado
para planejamento e despacho hidrotérmico em diversos países. À medida que
os modelos se tornam mais detalhados, em resolução horária, e a quantidade
de cenários aumenta, a execução do SDDP passa a demandar não apenas alto
poder computacional, mas também operações intensivas de E/S~\cite{MvLm:91}.

Nesse contexto, esta dissertação investiga o desempenho do MPI-IO sobre
sistemas de arquivos paralelos como o Lustre nessa classe de aplicações,
com foco no comportamento da E/S em ambientes HPC.

Esta dissertação concentra-se na fase de simulação final do SDDP ---
etapa executada após a convergência da política de operação, na qual um
conjunto fixo de cortes de Benders é avaliado em larga escala sobre
múltiplos cenários para produzir os resultados operacionais entregues ao
planejamento. As escolhas algorítmicas relacionadas à construção da
política (fase iterativa de geração de cortes) estão fora do escopo
deste trabalho.

\textcolor{red}{\textbf{ToDp: Apresentar organização da dissertação ao final
da introdução (capítulos 2--5).}}

